%------------------------------------------------------------------------------%

\usepackage{calculo_varias_variaveis-1}

\title {Derivação Implicita}

%------------------------------------------------------------------------------%
\begin{document}

\addtitle

%------------------------------------------------------------------------------%
\section{Derivação Implicita}

%------------------------------------------------------------------------------%
\begin{frame}{Derivação Implicita}
  Suponha que $F(x, y)$ seja diferenciável e que a equação $F(x, y)=0$
  defina $y$ como função de $x$.
\end{frame}

%------------------------------------------------------------------------------%
\begin{frame}{Derivação Implicita -- Justificativa}
\begin{columns}
\column{0.5\linewidth}
\begin{align*}
  \onslide<+->{w(x)          & = F\big(x, y(x)\big) = 0                              \\[2mm]}
  \onslide<+->{\frac{dw}{dx} & = \D{F}{x} \frac{dx}{dx} + \D{F}{y} \frac{dy}{dx} = 0 \\[2mm]}
  \onslide<+->{\frac{dw}{dx} & = F_x \frac{dx}{dx} + F_y \frac{dy}{dx} = 0}
\end{align*}
\column{0.5\linewidth}
\begin{align*}
  \onslide<+->{F_x + F_y \frac{dy}{dx} & = 0                 &  \\[2mm]}
  \onslide<+->{F_y \frac{dy}{dx}       & = - F_x             &  \\[2mm]}
  \onslide<+->{\frac{dy}{dx}           & = - \frac{F_x}{F_y} &  & \text{desde que } F_y \neq 0}
\end{align*}
\end{columns}
\end{frame}

%------------------------------------------------------------------------------%
\begin{frame}{Derivação Implicita -- Fórmula}
  Suponha que $F(x, y)$ seja diferenciável e que a equação $F(x, y)=0$
  defina $y$ como função de $x$

  \pause
  Então
  \[
    \frac{dy}{dx} = -\frac{F_x}{F_y}
  \]
\end{frame}

%------------------------------------------------------------------------------%
\addtocounter{exemplo}{1}
\begin{frame}{Exemplo \theexemplo}
  Calcule \(\frac{dy}{dx}\) se \(y^2-x^2-\sin(xy) = 0\)
\end{frame}

\begin{frame}{Exemplo \theexemplo}
  \[
    \frac{dy}{dx} = -\frac{F_x}{F_y}
  \]
  \pause
  \[
    F_x
    \pause
    = \D{}{x}\big(y^2-x^2-\sin(xy)\big)
    \pause
    = 0 - 2x - y\cos(xy)
    \pause
    = - 2x - y\cos(xy)
  \]
  \pause
  \[
    F_y
    \pause
    = \D{}{y}\big(y^2-x^2-\sin(xy)\big)
    \pause
    = 2y - 0 - x\cos(xy)
    \pause
    = 2y - x\cos(xy)
  \]
  \pause
  \[
    \frac{dy}{dx}
    \pause
    = -\frac{F_x}{F_y}
    \pause
    = -\frac{- 2x - y\cos(xy)}{2y - x\cos(xy)}
    \pause
    = \frac{2x + y\cos(xy)}{2y - x\cos(xy)}
  \]
\end{frame}

%------------------------------------------------------------------------------%
\begin{frame}{Derivação Implicita -- Três Variáveis}
  Suponha que $F(x, y, z)$ seja diferenciável e que a equação $F(x, y, z)=0$
  defina $z$ como função de $z=f(x, y)$
\end{frame}

%------------------------------------------------------------------------------%
\begin{frame}{Derivação Implicita -- Justificativa}
\begin{columns}
  \column{0.5\linewidth}
  \begin{align*}
    \onslide<+->{F\big(x, y, z(x, y)\big)                               & = 0 \\[2mm]}
    \onslide<+->{\D{}{x}\left[F\big(x, y, z(x, y)\big)\right]           & = 0 \\[2mm]}
    \onslide<+->{\D{F}{x}\D{x}{x} + \D{F}{y}\D{y}{x} + \D{F}{z}\D{z}{x} & = 0 \\[2mm]}
    \onslide<+->{\D{F}{x} + 0 + \D{F}{z}\D{z}{x} & = 0}
  \end{align*}
  \column{0.5\linewidth}
  \begin{align*}
    \onslide<+->{F_x + F_z \D{z}{x} & = 0                 &  &  \\[2mm]}
    \onslide<+->{F_z \D{z}{x}       & = - F_x             &  &  \\[2mm]}
    \onslide<+->{\D{z}{x}           & = - \frac{F_x}{F_z} &  & \text{desde que } F_z \neq 0}
  \end{align*}
\end{columns}
\end{frame}

%------------------------------------------------------------------------------%
\begin{frame}{Derivação Implicita -- Três Variáveis -- Fórmulas}
  Suponha que $F(x, y, z)$ seja diferenciável e que a equação $F(x, y, z)=0$
  defina $z$ como função de $z=f(x, y)$
  \[
    \D{z}{x} = -\frac{F_x}{F_z} \qquad
    \D{z}{y} = -\frac{F_y}{F_z}
  \]
\end{frame}

%------------------------------------------------------------------------------%
\addtocounter{exemplo}{1}
\begin{frame}{Exemplo \theexemplo}
  Calcule \(\D{z}{x}\) e \(\D{z}{z}\) se \(x^3 + z^2 + ye^{xz} + z\cos(y) = 0\)
\end{frame}

\begin{frame}{Exemplo \theexemplo}
  \[
    F(x, y, z) = x^3 + z^2 + ye^{xz} + z\cos(y)
  \]
  \pause
  \[
    F_x
    \pause
    = \D{}{x}\left[x^3 + z^2 + ye^{xz} + z\cos(y)\right]
    \pause
    = 3x^2 + 0 + ye^{xz}z + 0
    \pause
    = 3x^2 + yze^{xz}
  \]
  \pause
  \[
    F_y
    \pause
    = \D{}{y}\left[x^3 + z^2 + ye^{xz} + z\cos(y)\right]
    \pause
    = 0 + 0 + e^{xz} -z\sin(y)
    \pause
    = e^{xz} -z\sin(y)
  \]
  \pause
  \[
    F_z
    \pause
    = \D{}{z}\left[x^3 + z^2 + ye^{xz} + z\cos(y)\right]
    \pause
    = 0 + 2z + ye^{xy}x + \cos(y)
    \pause
    = 2z + yxe^{xy} + \cos(y)
  \]
\end{frame}

\begin{frame}{Exemplo \theexemplo}
  \[
    \D{z}{x}
    = -\frac{F_x}{F_z}
    \pause
    = -\frac{3x^2 + yze^{xz}}{2z + yxe^{xy} + \cos(y)}
  \]
  \pause
  \[
    \D{z}{y}
    = -\frac{F_y}{F_z}
    \pause
    = -\frac{e^{xz} -z\sin(y)}{2z + yxe^{xy} + \cos(y)}
  \]
\end{frame}

%------------------------------------------------------------------------------%
\section{Lista Mínima}

%------------------------------------------------------------------------------%
\begin{frame}{Lista Mínima}
  Cálculo Vol. 2 do Thomas 12\textsuperscript{a} ed. -- Seção 14.4
  \begin{enumerate}
    \item Estudar o texto da seção
    \item Resolver os exercícios: 26, 28, 30, 32
  \end{enumerate}

  \vfill
  Atenção: A prova é baseada no livro, não nas apresentações
\end{frame}

%------------------------------------------------------------------------------%
\end{document}
%------------------------------------------------------------------------------%
